\documentclass[11pt]{article}
\usepackage[margin=1in]{geometry}
\usepackage{booktabs}
\usepackage{amsmath}
\usepackage{url}
\usepackage{hyperref}
\usepackage{graphicx}

\title{Representation Is the Variable:\\
A Machine-Verified Benchmark for Functional Harmony}
\author{Axel Stevens\\ \texttt{axel@aesv.io}}
\date{}

\begin{document}
\maketitle

\begin{abstract}
Language models answer music-theory questions well and analyse music badly, and
existing benchmarks cannot separate the two because they vary the question and
the encoding at the same time. We present a benchmark of 21{,}940 items in which
the same harmonic content is presented in four encodings of increasing
difficulty --- chord symbols, spelled notes, and bass-first pitch-class sets for
Roman-numeral analysis, plus a key-identification framing --- so that
representation is the only variable. Gold labels are never hand-annotated: a
probabilistic common-practice grammar emits each progression together with its
intended analysis, and every chord is then re-derived along two independent
\texttt{music21} paths, with an item retained only if both agree on pitch-class
set and bass. All 27{,}480 chords agree. Because no annotated corpus is
redistributed, the benchmark carries a CC-BY-4.0 licence, unlike the
non-commercial corpora that dominate this task. Scoring is exact match against
deterministic gold, with no model judge. Evaluating six models zero-shot, we find
that no model saturates even the easiest framing: the best exact match on chord
symbols is 0.256, and models without inference-time reasoning collapse toward
zero once chord symbols are withheld, while per-chord accuracy degrades far more
gracefully than sequence-level exact match. A paired ablation shows the benchmark
measures different capabilities depending on the decoding regime: with reasoning
disabled it measures pattern recall, and with reasoning enabled it measures
multi-step computation, with exact match rising from 0.000 to 0.740 for one model
on the same items. Neither regime saturates it.
\end{abstract}

\section{Introduction}

Roman-numeral analysis asks, of each chord in a passage, what function it serves
relative to a tonic. It is the central skill of tonal analysis and the first one
taught, and unlike most music-understanding tasks it has a largely
agreed-upon symbolic answer. That makes it an unusually clean target for
evaluating language models: the answer is a short string, it is checkable
without a judge, and getting it right requires composing several steps ---
identify the notes, stack them into a chord, locate the chord's root relative to
the key, name the inversion, and decide what the final motion cadentially is.

Language models appear to know a great deal of music theory when asked about it
in prose. Whether they can perform the analysis is a separate question, and it is
the question existing benchmarks answer least well, because they conflate two
things. A multiple-choice item about the definition of a half cadence and an
open-ended request to analyse a progression differ both in what is being asked
and in how the music is encoded. If a model does well on the first and poorly on
the second, the result is not attributable.

We isolate the encoding. Every item in our benchmark is drawn from the same pool
of progressions, and each progression is presented in more than one form:

\begin{itemize}
  \item \textbf{\texttt{symbol\_to\_rn}} gives the key and printed chord symbols
    (\texttt{Dm7 G7 Cmaj7}). Chord quality is handed to the model; the remaining
    work is functional.
  \item \textbf{\texttt{notes\_to\_rn}} gives the key and spelled notes
    (\texttt{D4 F4 A4 C5 | \dots}). The model must read each simultaneity as a
    chord before it can name a function.
  \item \textbf{\texttt{pcset\_to\_rn}} gives the key and bass-first pitch-class
    integers (\texttt{[2 5 9 0] | \dots}). Spelling is gone; enharmonic
    distinctions must be recovered from context.
  \item \textbf{\texttt{key\_id}} withholds the key entirely and asks only for it,
    restricted to progressions whose key is uniquely determined.
\end{itemize}

The harmonic content is fixed across the three Roman-numeral framings. Any
difference in score between them is a difference in the model's ability to
recover harmony from a representation, not a difference in the underlying
question. Our contributions are:

\begin{enumerate}
  \item A benchmark of 21{,}940 items in which representation is a controlled
    variable, built from 845 key-independent progression shapes.
  \item A gold-label pipeline with no human annotation and a verifiable
    agreement gate: dual derivation through two independent \texttt{music21}
    entry points, 27{,}480 of 27{,}480 chords in agreement.
  \item A permissive licence for a task whose standard corpora are
    non-commercial, achieved by generating rather than redistributing.
  \item Zero-shot results for six models and a paired reasoning-on/off ablation
    showing that the benchmark measures a different capability in each regime.
\end{enumerate}

\section{Related work}

\textbf{Music-theory question answering.} MusicTheoryBench \cite{chatmusician}
contributes 372 hand-written multiple-choice questions spanning broad music
knowledge. It measures recall of theory facts across many topics; we measure
execution of one task across many encodings, and generate rather than
hand-write, which is what lets us scale to five figures and control the
representation.

\textbf{Synthetic music reasoning.} Kruspe \cite{kruspe} generates synthetic
interval, chord, and scale identification problems for language models. That
work shares our generate-don't-annotate stance but supplies no key context, so
the task is identification rather than functional analysis: a chord is named in
isolation, not situated against a tonic.

\textbf{Encoding as a variable.} Pond and Fujinaga \cite{pond} present a single
RCM Level 6 examination in four encodings together with a study of prompting
strategies, and are the closest prior work in spirit. Their design varies
encoding across a small fixed instrument; ours varies encoding across thousands
of generated items drawn from one pool, which trades their ecological validity
for statistical power and a controlled comparison.

\textbf{Specialist analysis models.} A substantial line of work trains dedicated
models for score-based Roman-numeral analysis, including AugmentedNet
\cite{augmentednet}, graph neural approaches \cite{rnagnn}, and AnalysisGNN
\cite{analysisgnn}. These systems target scores rather than text prompts and
train on human-annotated corpora --- ABC, BPS, HaydnSun, TAVERN, When-in-Rome
\cite{wir}, WTC --- most of which are licensed non-commercially. Dilemmadata
\cite{dilemmadata} documents the interoperability problems that follow from
combining them. Our benchmark is complementary: it does not compete with
specialist accuracy on real repertoire, and it is not trained on. It measures
what general-purpose text models can do on harmony they have never been fitted
to, under a licence that permits unrestricted use.

\section{Construction}

\subsection{Grammar}

Progressions are emitted by a probabilistic grammar over common-practice
functional syntax: a tonic region, an optional predominant, a dominant, and a
cadential resolution, with sevenths, inversions, cadential six-four figures, and
secondary dominants available at each step. The grammar emits a
\emph{shape}: a key-independent Roman-numeral sequence together with its intended
cadence label. Shapes are then transposed across keys to produce instances. This
separation matters for the splits (Section~\ref{sec:splits}).

The label convention follows the feature decomposition of the DCML harmony
standard \cite{dcml} --- \texttt{numeral}, \texttt{form}, \texttt{figbass},
\texttt{changes}, \texttt{relativeroot}. We adopt the notation and copy no DCML
data. Cadences are coded \texttt{PAC}, \texttt{IAC}, \texttt{HC}, \texttt{DC},
\texttt{PC}.

\subsection{Verification by dual derivation}

The grammar knows what it intended, but an intention is not a label. Every chord
is independently re-derived along two paths that share no code:

\begin{enumerate}
  \item the Roman-numeral figure is realised through the \texttt{music21}
    \cite{music21} Roman-numeral engine;
  \item the printed chord symbol is parsed through the \texttt{music21}
    chord-symbol parser.
\end{enumerate}

An item is retained only if both paths agree on pitch-class set and bass. This
catches a class of error that a single derivation cannot: a figure that the
grammar composed incorrectly will usually still round-trip through one engine,
but will disagree with the other. On the released build, 27{,}480 of 27{,}480
chords agree and zero instances were dropped. We report this as a gate rather
than a result: it is a necessary condition, not evidence that the harmonic
distribution is musically natural.

\subsection{Splits}
\label{sec:splits}

The atomic unit is the shape, not the item. A shape hashes to exactly one split,
so no transposition of a progression --- and no alternative framing of it ---
can appear on both sides of the train/test boundary. Without this, a model could
see \texttt{ii7 V7 IM7} in C at training time and be tested on the same shape in
E$\flat$, or on the same progression rendered as pitch classes. Rows divide
14{,}725 / 3{,}571 / 3{,}644 across train, validation, and test. The test split
is the benchmark.

\subsection{Key-identification gate}

Gold keys are not always well defined. In C major, \texttt{I V7/V V} is
note-identical to \texttt{IV V7 I} in G major, and the G reading is arguably the
stronger one. Rather than adjudicate, we exclude: \texttt{key\_id} retains only
cadence-terminated progressions of three or more chords whose notes contain both
scale degree 4 and the leading tone, which pins the key uniquely. This costs
coverage and buys defensible gold.

\section{Evaluation protocol}

Gold is deterministic, so scoring is exact string match after normalisation. No
model judge is involved. For the Roman-numeral configurations we report
\texttt{exact} (numerals and cadence both correct), \texttt{labels\_exact}
(numerals correct, cadence ignored), \texttt{chord\_acc} (per-chord accuracy),
and \texttt{cadence\_acc}. For \texttt{key\_id} we report \texttt{exact},
\texttt{tonic\_acc}, and \texttt{mode\_acc}. Prompts are versioned; all results
below use prompt v1, zero-shot, temperature 0. The harness depends only on the
Python standard library and speaks to any OpenAI-compatible endpoint.

We also report \texttt{parse\_failures}: responses from which no answer could be
extracted. These are counted as incorrect, but reported separately, because a
model that analyses correctly and formats disobediently is failing at something
other than harmony, and the two should not be silently merged.

\section{Results}

\subsection{Zero-shot}

Table~\ref{tab:main} reports the full test split, run 2026-07-11 via OpenRouter
at a total API cost of approximately nine US dollars. Cells give exact match with
per-chord accuracy in parentheses.

\begin{table}[h]
\centering
\small
\begin{tabular}{lcccc}
\toprule
model & \texttt{symbol\_to\_rn} & \texttt{notes\_to\_rn} & \texttt{pcset\_to\_rn} & \texttt{key\_id} \\
 & $n{=}943$ & $n{=}943$ & $n{=}943$ & $n{=}815$ \\
\midrule
gpt-oss-120b (reasoning low) & \textbf{0.256} (0.885) & \textbf{0.155} (0.778) & 0.146 (0.724) & 0.778 \\
Claude Sonnet 5              & 0.220 (0.863) & 0.066 (0.688) & \textbf{0.157} (0.732) & \textbf{0.823} \\
Qwen3-235B-A22B-Instruct     & 0.193 (0.739) & 0.016 (0.380) & 0.002 (0.163) & 0.719 \\
DeepSeek-V3.2                & 0.151 (0.634) & 0.005 (0.305) & 0.001 (0.150) & 0.661 \\
Kimi-K2.5                    & 0.142 (0.562) & 0.007 (0.399) & 0.009 (0.205) & 0.788 \\
Llama-3.3-70B                & 0.037 (0.466) & 0.004 (0.307) & 0.000 (0.194) & 0.471 \\
\bottomrule
\end{tabular}
\caption{Zero-shot, temperature 0, prompt v1. Exact match, per-chord accuracy in
parentheses. Kimi-K2.5 incurred 156, 28, and 141 parse failures on the three
Roman-numeral configurations respectively; its scores understate its harmonic
ability by an unquantified margin.}
\label{tab:main}
\end{table}

Three observations. First, no model saturates the easiest configuration: the
best exact match where chord quality is given outright is 0.256. Second, the drop
from \texttt{symbol\_to\_rn} to \texttt{notes\_to\_rn} is severe for models that
do not reason at inference time --- Qwen3, DeepSeek, Kimi, and Llama all fall
below 0.02 --- while the two that do degrade far more slowly. This is the
signature of chord-symbol lookup: a model that recognises \texttt{Dm7} as a token
and retrieves \texttt{ii7} has nothing to retrieve once the token is gone.
Third, the gap between exact match and per-chord accuracy is large everywhere.
Llama-3.3-70B scores 0.004 exact on \texttt{notes\_to\_rn} while getting 30.75\%
of individual chords right; sequence-level exact match is an unforgiving metric
and per-chord accuracy is the better instrument for tracking partial competence.

Claude Sonnet 5 is the one model that scores \emph{higher} on pitch classes
(0.157) than on spelled notes (0.066). We do not have a mechanism for this and
flag it as an anomaly worth investigating; the natural hypothesis is that
note-spelling parsing, rather than harmonic reasoning, is its bottleneck on
\texttt{notes\_to\_rn}, but we have not tested it.

\subsection{Reasoning on versus off}

Table~\ref{tab:reasoning} reports paired runs on fixed test subsets: identical
prompts and items, reasoning disabled and enabled.

\begin{table}[h]
\centering
\small
\begin{tabular}{llcccccccc}
\toprule
& & \multicolumn{2}{c}{\texttt{symbol\_to\_rn}} & \multicolumn{2}{c}{\texttt{notes\_to\_rn}} & \multicolumn{2}{c}{\texttt{pcset\_to\_rn}} & \multicolumn{2}{c}{\texttt{key\_id}} \\
\cmidrule(lr){3-4}\cmidrule(lr){5-6}\cmidrule(lr){7-8}\cmidrule(lr){9-10}
model & $n$ & off & on & off & on & off & on & off & on \\
\midrule
Kimi-K2.5     & 150 & 0.107 & 0.907 & 0.000 & 0.740 & 0.013 & 0.840 & 0.760 & 0.913 \\
DeepSeek-V3.2 & 200 & 0.145 & 0.805 & 0.025 & 0.725 & 0.005 & 0.770 & 0.630 & 0.940 \\
Claude Sonnet 5 & 100 & 0.220 & 0.590 & 0.090 & 0.520 & 0.140 & 0.660 & 0.860 & 0.960 \\
gpt-oss-120b (low$\to$high) & 200 & 0.255 & 0.580 & 0.185 & 0.440 & 0.120 & 0.535 & 0.785 & 0.855 \\
\bottomrule
\end{tabular}
\caption{Exact match, reasoning disabled versus enabled, on fixed subsets of the
test split. Same prompts, same items.}
\label{tab:reasoning}
\end{table}

The effect is large and consistent across every model and configuration. Kimi-K2.5
moves from 0.000 to 0.740 on \texttt{notes\_to\_rn}; DeepSeek-V3.2 from 0.005 to
0.770 on \texttt{pcset\_to\_rn}. With reasoning enabled, per-chord accuracy on the
harder configurations reaches 0.89--0.98, so most of the remaining exact-match
gap is cadence labelling and figured-bass detail rather than chord identification.

We read this as a statement about what the benchmark measures rather than about
which models are better. Without inference-time computation, the task is
answerable only by recall, and the scores reflect how much of this material a
model has memorised. With it, the task becomes the multi-step derivation it was
designed to be. A benchmark that behaves this differently under the two regimes
should always be reported with the regime attached.

\section{Limitations}

\textbf{The distribution is synthetic.} These are grammar outputs, not
repertoire. Chord-transition statistics are those of the grammar, not of any
composer, and results here do not license claims about performance on real
scores. This is the benchmark's central limitation and no amount of verification
addresses it.

\textbf{Parse failures confound.} Kimi-K2.5's 156 parse failures on
\texttt{symbol\_to\_rn} mean its 0.142 conflates format non-compliance with
harmonic error. Every table in this paper counts unparseable output as wrong. A
prompt-robustness study would separate them and we have not run one.

\textbf{Single prompt, single run.} All results use prompt v1 at temperature 0,
one run per cell. We report no variance and have not measured prompt
sensitivity, which is known to be substantial for tasks of this shape.

\textbf{Small ablation subsets.} The reasoning comparison uses 100--200 items per
cell, chosen for cost. The effect sizes are large enough that the direction is
not in question, but the point estimates are loose.

\textbf{One planned comparison did not run.} A Qwen3-235B Thinking versus
Instruct checkpoint pair was scoped and is absent from the results.

\textbf{Cadence rules are strict.} PAC versus IAC is decided by inversion, since
no soprano is notated. A half cadence must end on a root-position \texttt{V}
triad, so a terminal \texttt{V7} is unlabelled; a deceptive cadence requires a
root-position dominant, so \texttt{V65}$\to$\texttt{vi} does not count. These are
defensible conventions, not the only ones, and a model applying a different
convention consistently will be scored as wrong.

\textbf{Harmony only.} No voice leading, melody, or rhythm. No modal mixture,
Neapolitan sixths, or augmented sixths.

\section{Conclusion}

Holding harmonic content fixed while varying its encoding turns out to separate
two things that benchmarks in this area normally merge: knowing music theory and
doing it. Models that cannot do it still score respectably when chord symbols are
supplied and collapse when they are withheld, and the collapse is a property of
the decoding regime as much as of the model. The benchmark, its generator, its
harness, and its verification report are released under CC-BY-4.0, and because
nothing annotated is redistributed, they can be used without the licence
friction that attaches to the standard corpora for this task.

\begin{thebibliography}{99}
\bibitem{chatmusician} R. Yuan et al. ChatMusician: Understanding and Generating
Music Intrinsically with LLM. 2024. MusicTheoryBench.
\bibitem{kruspe} A. Kruspe. Harmonic Reasoning in Large Language Models.
arXiv:2409.05521, 2024.
\bibitem{pond} T. Pond and I. Fujinaga. Teaching Large Language Models Music
Theory. arXiv:2503.22853, 2025.
\bibitem{augmentednet} G. Micchi, M. Gotham, M. Giraud. AugmentedNet: A Roman
Numeral Analysis Network with Synthetic Training Examples and Additional Tonal
Tasks. ISMIR, 2020.
\bibitem{rnagnn} E. Karystinaios and G. Widmer. Roman Numeral Analysis with Graph
Neural Networks. ISMIR, 2023.
\bibitem{analysisgnn} AnalysisGNN: Unified Music Analysis with Graph Neural
Networks. arXiv:2509.06654, 2025.
\bibitem{wir} M. Gotham et al. When in Rome: a meta-corpus for functional
harmonic analysis. \url{https://github.com/MarkGotham/When-in-Rome}
\bibitem{dilemmadata} Dilemmadata: On the Interoperability of Heterogeneous Roman
Numeral Datasets. arXiv:2606.31595, 2026.
\bibitem{dcml} DCML harmony annotation standard.
\url{https://github.com/DCMLab/standards}
\bibitem{music21} M. Cuthbert and C. Ariza. music21: A Toolkit for
Computer-Aided Musicology and Symbolic Music Data. ISMIR, 2010.
\end{thebibliography}

\end{document}
